% GAL1897_W09: modern English translation, PDF 062--079 / L0061--L0078.
% The claims and defects of the 1897 printing are preserved; corrections and
% witness history appear only in explicitly identified critical notes.

\providecommand{\GalSourcePageStart}[4]{}%
\providecommand{\GalSourcePageBoundary}[4]{\space}%
\providecommand{\GalPrintedError}[2]{#1}%
\providecommand{\GalPrintedOmission}[1]{}%
\providecommand{\GalSourceError}[2]{#1}%
\providecommand{\GalWitnessVariant}[2]{#1}%
\providecommand{\GalControlReading}[2]{#1}%
\providecommand{\GalSourceErrorMark}[1]{}%
\providecommand{\GalProofGapMark}[1]{}%
\providecommand{\GalControlDefect}[2]{#1}%
\providecommand{\GalHistoricalNote}[1]{\par\smallskip{\footnotesize #1\par}}%
\providecommand{\GalArticleEndRule}{\par\smallskip\begin{center}\rule{0.18\linewidth}{0.4pt}\end{center}}%

\GalSourcePageStart{062}{L0061}{33}{33}

% ENSEG:W09:PDF062:0001
\begin{center}
{\Large\bfseries MEMOIR\par}
\vspace{0.35em}
{\small\bfseries ON THE\par}
\vspace{0.35em}
{\large\bfseries CONDITIONS FOR THE SOLVABILITY OF EQUATIONS BY RADICALS\textsuperscript{(*)}.\par}
\vspace{0.8em}
\rule{0.12\linewidth}{0.4pt}
\end{center}

% ENSEG:W09:PDF062:0002
The accompanying Memoir\textsuperscript{(2)} is extracted from a work that I had
the honor of presenting to the Academy a year ago. Since that work was not
understood and the propositions it contains were called into question, I have
had to content myself with giving, in synthetic form, the general principles
and a \emph{single} application of my theory. I beg my judges at least to read
these few pages attentively.

% ENSEG:W09:PDF062:0003
Here one will find a general \emph{condition satisfied by every equation solvable
by radicals}, and which, conversely, guarantees solvability. It is applied only to

% ENSEG:W09:PDF062:0004
\GalHistoricalNote{\textsuperscript{(*)} This Memoir and the one following it were found among Galois's papers and
were first published in 1846 by Liouville, who prefaced them with the following
note:

``When they inserted in their periodical the letter just read, the editors of the
\emph{Revue encyclopédique} announced that they would soon publish the manuscripts
left by Galois. But \GalControlDefect{this}{W09-CD001} promise was not kept. M.~Auguste Chevalier
had nevertheless prepared the work. He has given it to us, and in the sheets that
follow one will find:

``1\textsuperscript{o} A complete Memoir on the conditions for the solvability of equations by
radicals, with the application to equations of prime degree;

``2\textsuperscript{o} A fragment of a second Memoir in which Galois treats the general theory
of equations that he calls \emph{primitive}.

``We have retained most of the notes that M.~Auguste Chevalier appended to the
Memoirs just mentioned. These notes are all marked with the initials A.~Ch. The
unsigned notes are by Galois himself.

``We shall complete this publication with a few other pieces extracted from
Galois's papers which, though not of great importance, may nevertheless still
be read with interest by geometers.''

The extracts mentioned by Liouville in the last sentence of this note were never
published.

\textsuperscript{(2)} I thought it appropriate to place at the head of this Memoir the preface about to
be read, although I found it crossed out in the manuscript. \hfill \textsc{(A.~Ch.)}

\hfill E.~G.}

\GalSourcePageBoundary{063}{L0062}{34}{34}

% ENSEG:W09:PDF063:0005
equations whose degree is a prime number. Here is the theorem yielded by our
analysis:

% ENSEG:W09:PDF063:0006
\begin{quote}\itshape
For an equation of prime degree having no commensurable divisors to be
solvable by radicals, it is necessary and sufficient that all its roots be
rational functions of any two of them.
\end{quote}

% ENSEG:W09:PDF063:0007
The other applications of the theory each constitute a separate special
theory. They also require the use of number theory and of a particular
algorithm; we reserve them for another occasion. They concern in part the
modular equations of the theory of elliptic functions, which we prove cannot
be solved by radicals.

% ENSEG:W09:PDF063:0008
\medskip
\noindent\hspace{3em}16 January 1831.\hfill \textsc{E. Galois.}\hspace{3em}

% ENSEG:W09:PDF063:0009
\begin{center}
\rule{0.15\linewidth}{0.4pt}\par
\vspace{0.7em}
{\bfseries PRINCIPLES.}
\end{center}

% ENSEG:W09:PDF063:0010
I shall begin by establishing a few definitions and a sequence of lemmas, all
of which are known.

% ENSEG:W09:PDF063:0011
\emph{Definitions.} --- An equation is called \emph{reducible} when it has rational
divisors, and \emph{irreducible} otherwise.

% ENSEG:W09:PDF063:0012
Here it is necessary to explain what is meant by the word \emph{rational}, since
it will occur often.

% ENSEG:W09:PDF063:0013
When the equation has \emph{all} its coefficients numerical and rational, this
means simply that the equation can be decomposed into factors whose
coefficients are numerical and rational.

% ENSEG:W09:PDF063:0014
But when the coefficients of an equation are \emph{not all} numerical and
rational, a rational divisor must then be understood to mean a divisor whose
coefficients can be expressed as rational functions of the coefficients of the
given equation.\GalPrintedOmission{W09-WV001}

% ENSEG:W09:PDF063:0015
\GalCriticalNote{POST-P13-A014}{\textbf{Status: proved; the W09 provenance corpus is closed to the 2026-08-13 search cutoff.}
The memoir was collated against the 1846 publication, Tannery 1908, the
1962/1976 critical-edition records, Neumann 2011/2013, and Dicker 2026. Earlier
witnesses, explicit notices, later reconstructions, and bounded no-notice
findings remain separately identified in the returned 51-row priority matrix.}

\GalCriticalNote{W09-WV001}{\textbf{Status: documented 1897 omission; explicit earlier and later witnesses located.}
After the sentence on rational divisors, the 1897 printing omits two definitional
lines that appear in the 1846 publication and the 1908 collation and give a
general definition of a ``rational quantity.'' They are not restored here. The
omission makes the transition less explicit but does not authorize importing
another witness into the translated 1897 text. The 1846 publication supplies
the lines, and Tannery explicitly records the omission in 1908. Evidence:
W09-E063, W09-E1846-039, W09-E1908-014, and POST-P13-A014.}

% ENSEG:W09:PDF063:0016
More generally, one may agree to regard as rational every rational function of
a certain number of specified quantities assumed to be known \emph{a priori}.
For example, one may

\GalSourcePageBoundary{064}{L0063}{35}{35}

% ENSEG:W09:PDF064:0017
choose a particular root of an integer and regard every rational function of
that radical as rational.

% ENSEG:W09:PDF064:0018
When we agree in this way to regard certain quantities as known, we shall say
that we \emph{adjoin} them to the equation to be solved. We shall say that these
quantities are \emph{adjoined} to the equation.

% ENSEG:W09:PDF064:0019
This being established, we shall call \emph{rational} every quantity that can be
expressed as a rational function of the coefficients of the equation and of a
certain number of quantities \emph{adjoined} to the equation and chosen
arbitrarily.

% ENSEG:W09:PDF064:0020
When we use auxiliary equations, they will be rational if their coefficients
are rational in our sense.

% ENSEG:W09:PDF064:0021
It is clear, moreover, that the properties and difficulties of an equation may
be entirely different according to the quantities adjoined to it. For example,
adjoining one quantity may make an irreducible equation reducible.

% ENSEG:W09:PDF064:0022
Thus, when one adjoins to the equation

% ENSEG:W09:PDF064:0023
\[
\frac{x^{n}-1}{x-1}=0\GalControlReading{.}{W09-CD007}
\]

% ENSEG:W09:PDF064:0024
where $n$ is prime, a root of one of M.~Gauss's auxiliary equations, the
equation decomposes into factors and consequently becomes reducible.

% ENSEG:W09:PDF064:0025
Substitutions are the passage from one permutation to another.

% ENSEG:W09:PDF064:0026
The permutation from which one starts in order to specify the substitutions is
entirely arbitrary when functions are concerned, for there is no reason that,
in a function of several letters, one letter should occupy one position rather
than another.

% ENSEG:W09:PDF064:0027
Nevertheless, because it is difficult to form the idea of a substitution
without also forming that of a permutation, we shall frequently use
permutations in our terminology and shall regard substitutions only as the
passage from one permutation to another.

% ENSEG:W09:PDF064:0028
When we wish to group substitutions, we shall make them all arise from the same
permutation.

% ENSEG:W09:PDF064:0029
Since the original arrangement of the letters has no effect in the questions
concerning the groups that we shall consider, the substitutions must be the
same whatever

\GalSourcePageBoundary{065}{L0064}{36}{36}

% ENSEG:W09:PDF065:0030
the permutation from which one starts. Thus, if such a group contains the
substitutions $S$ and $T$, it certainly contains the substitution $ST$.

% ENSEG:W09:PDF065:0031
These are the definitions we thought it necessary to recall.

% ENSEG:W09:PDF065:0032
\medskip
\noindent\textsc{Lemma I.} --- \emph{An irreducible equation cannot have a root in
common with a rational equation without dividing it.}

% ENSEG:W09:PDF065:0033
For the greatest common divisor of the irreducible equation and the other
equation will itself be rational; hence, etc.

% ENSEG:W09:PDF065:0034
\medskip
\noindent\textsc{Lemma II.} --- \emph{Given any equation without repeated roots,
whose roots are $a,b,c,\ldots$, one can always form a function $V$ of the roots
such that no two values obtained by permuting the roots in
\GalPrintedError{that}{W09-PE001} function in every possible way are equal.}

% ENSEG:W09:PDF065:0035
For example, one may take

% ENSEG:W09:PDF065:0036
\[
V=Aa+Bb+Cc+\ldots,
\]

% ENSEG:W09:PDF065:0037
where $A,B,C$ are suitably chosen integers.

% ENSEG:W09:PDF065:0038
\GalCriticalNote{W09-PE001}{\textbf{Status: proved printed-word error.}
The 1897 French reads the ill-formed \emph{celle fonction}; the English main
text preserves its demonstrative as ``that function.'' The supported reading
is \emph{cette fonction}, ``this function,'' which the 1846 publication also
prints. The correction changes no mathematical claim and is confined to this
note. The 1846 witness supplies the correct form, but the bounded search found
no explicit itemized notice. Evidence: W09-E065, W09-E1846-040,
POST-P13-A009, and POST-P13-A014.}

\GalSourceErrorMark{W09-SE001}
% ENSEG:W09:PDF065:0039
\medskip
\noindent\textsc{Lemma III.} --- \emph{When the function $V$ has been chosen as
indicated in the preceding article, it will have the property that all the
roots of the given equation can be expressed rationally as functions of $V$.}

% ENSEG:W09:PDF065:0040
Indeed, let

% ENSEG:W09:PDF065:0041
\[
V=\varphi(a,b,c,d,\ldots),
\]

% ENSEG:W09:PDF065:0042
or

% ENSEG:W09:PDF065:0043
\[
V-\varphi(a,b,c,d,\ldots)=0.
\]

% ENSEG:W09:PDF065:0044
Multiply together all similar equations obtained by permuting all the letters
in these equations while leaving only the first fixed; the following
expression results:

% ENSEG:W09:PDF065:0045
\[
\begin{split}
&[V-\varphi(a,b,c,d,\ldots)]
[V-\varphi(a,c,b,d,\ldots)]\\
&\qquad [V-\varphi(a,b,d,c,\ldots)]\ldots,
\end{split}
\]

% ENSEG:W09:PDF065:0046
It is symmetric in $b,c,d,\ldots$ and can therefore be written as a function of
$a$. We shall thus have an equation of the form

% ENSEG:W09:PDF065:0047
\[
F(V,a)=0.
\]

\GalSourcePageBoundary{066}{L0065}{37}{37}

% ENSEG:W09:PDF066:0048
I now say that the value of $a$ can be obtained from this equation. For this it
is enough to seek the common solution of this equation and the given equation.
That solution is their only common one, for one cannot have, for example,

% ENSEG:W09:PDF066:0049
\[
F(V,b)=0,
\]

% ENSEG:W09:PDF066:0050
with this equation having a common factor with the similar equation; otherwise
one of the functions $\varphi(a,\ldots)$ would equal one of the functions
$\varphi(b,\ldots)$, contrary to the hypothesis.

% ENSEG:W09:PDF066:0051
It follows that $a$ can be expressed as a rational function of $V$, and the
same is true of the other roots.

% ENSEG:W09:PDF066:0052
\GalCriticalNote{W09-SE001}{\textbf{Status: repaired; Proposition I and later rational-resolvent uses survive.}
The printed argument merely says to seek the unique common root and then
declares it rational in $V$. The missing step is to work in $K(V)[X]$, prove
that the relevant greatest common divisor is $X-a$, and use the Euclidean
algorithm and a Bézout identity to obtain $a\in K(V)$. This repairs Lemma III
and supports its use in Proposition I; it does not alter the translated 1897
proof. Tannery's 1908 collation explicitly reports the insufficiency. Evidence:
W09-E065, W09-E066, W09-E1908-014, W09-MATH001, and POST-P13-A010.}

% ENSEG:W09:PDF066:0053
\GalCriticalNote{W09-WV002}{\textbf{Status: documented witness history; not an 1897 error.}
The 1908 collation reports that Galois regarded this proof as insufficient and
that Poisson/Galois material in copy or proof was suppressed by Liouville. That
history contextualizes but neither replaces nor silently expands the 1897
proof. Downstream mathematical repair is recorded under W09-SE001. Tannery's
1908 collation is the explicit witness notice. Evidence: W09-E1908-014 and
POST-P13-A014.}

% ENSEG:W09:PDF066:0054
This proposition\textsuperscript{(*)} is cited without proof by Abel in the
posthumous Memoir on elliptic functions.

% ENSEG:W09:PDF066:0055
\medskip
\noindent\textsc{Lemma IV.} --- \emph{Suppose that the equation in $V$ has been
formed and one of its irreducible factors has been selected, so that $V$ is a
root of an irreducible equation. Let $V,V',V'',\ldots$ be the roots of this
irreducible equation. If $a=f(V)$ is one root of the given equation, then
$f(V')$ will likewise be a root of the given equation.}

% ENSEG:W09:PDF066:0056
Indeed, multiplying together all factors of the form
$V-\varphi(a,b,c,\ldots,d)$ obtained by making all possible permutations of the
letters gives a rational equation in $V$, which is necessarily divisible by
the equation in question; hence $V'$ must be obtained by exchanging the
letters in the function $V$. Let $F(V,a)=0$ be the equation obtained by
permuting in $V$ all the letters except the first. We shall then have
$F(V',b)=0$, where $b$ may equal $a$ but is certainly one of the roots of the
given equation. Just as the given equation together with $F(V,a)=0$ yielded
$a=f(V)$, so the given equation together with $F(V',b)=0$ yields $b=f(V')$.

% ENSEG:W09:PDF066:0057
\GalHistoricalNote{\textsuperscript{(*)} It is remarkable that this proposition implies that every equation depends on
an auxiliary equation whose roots are all rational functions of one another,
for the auxiliary equation in $V$ is such an equation.

This observation is, however, merely curious. An equation having this property
is not in general easier to solve than any other.}

\GalSourcePageBoundary{067}{L0066}{38}{38}

% ENSEG:W09:PDF067:0058
\begin{center}{\bfseries PROPOSITION I.}\end{center}

% ENSEG:W09:PDF067:0059
\noindent\textsc{Theorem.} --- {\itshape Let an equation be given whose $m$ roots are
$a,b,c,\ldots$. There will always be a group of permutations of the letters
$a,b,c,\ldots$ having the following property:\par

1\textsuperscript{o} Every function of the roots invariant\textsuperscript{(*)} under the substitutions
of this group is rationally known;\par

2\textsuperscript{o} Conversely, every function of the roots that is rationally
determinate is invariant under the substitutions.\par}

% ENSEG:W09:PDF067:0060
(For algebraic equations, this group is nothing other than the set of the
$1.2.3\ldots m$ possible permutations of the $m$ letters, since in that case
only symmetric functions are rationally determinate.)

% ENSEG:W09:PDF067:0061
(For the equation $\dfrac{x^{n}-1}{x-1}=0$, if one sets $a=r$, $b=r^{g}$,
$c=r^{g^{2}},\ldots$, where $g$ is a primitive root, the group of permutations
is simply the following:

% ENSEG:W09:PDF067:0062
\[
\begin{array}{c}
abcd\ldots k,\\
bcd\ldots ka,\\
cd\ldots kab,\\
\ldots\ldots\ldots\\
kabc\ldots i;
\end{array}
\]

% ENSEG:W09:PDF067:0063
in this special case the number of permutations equals the degree of the
equation, and the same would hold for equations whose roots are all rational
functions of one another.)

% ENSEG:W09:PDF067:0064
\emph{Proof.} --- Whatever equation is given, one can find a rational function
$V$ of the roots such that all

% ENSEG:W09:PDF067:0065
\GalHistoricalNote{\textsuperscript{(*)} Here we call invariant not only a function whose form is invariant under
substitutions of the roots among themselves, but also one whose numerical value
does not vary under those substitutions. For example, if $Fx=0$ is an equation,
$Fx$ is a function of the roots that varies under no permutation.

When we say that a function is rationally known, we mean that its numerical
value can be expressed as a rational function of the coefficients of the
equation and of the adjoined quantities.}

\GalSourcePageBoundary{068}{L0067}{39}{39}

% ENSEG:W09:PDF068:0066
the roots are rational functions of $V$. This being established, consider the
irreducible equation of which $V$ is a root (Lemmas III and IV). Let
$V,V',V'',\ldots,V^{(n-1)}$ be the roots of this equation.

% ENSEG:W09:PDF068:0067
Let $\varphi V,\varphi_{1}V,\varphi_{2}V,\ldots,\varphi_{m-1}V$ be the roots of
the given equation.

% ENSEG:W09:PDF068:0068
Write the following $n$ permutations of the roots:

% ENSEG:W09:PDF068:0069
\[
\begin{array}{c@{\quad}ccccc}
(V)&\varphi V&\varphi_{1}V&\varphi_{2}V&\ldots&\varphi_{m-1}V,\\
(V')&\varphi V'&\varphi_{1}V'&\varphi_{2}V'&\ldots&\varphi_{m-1}V',\\
(V'')&\varphi V''&\varphi_{1}V''&\varphi_{2}V''&\ldots&\varphi_{m-1}V'',\\
\ldots&\ldots&\ldots&\ldots&\ldots&\ldots,\\
(V^{(n-1)})&\varphi V^{(n-1)}&\varphi_{1}V^{(n-1)}&\varphi_{2}V^{(n-1)}&\ldots&\varphi_{m-1}V^{(n-1)}:
\end{array}
\]

% ENSEG:W09:PDF068:0070
I say that this group of permutations has the stated property.

% ENSEG:W09:PDF068:0071
Indeed:

% ENSEG:W09:PDF068:0072
1\textsuperscript{o} Every function $F$ of the roots invariant under the
substitutions of this group can be written $F=\psi V$, and we shall have

% ENSEG:W09:PDF068:0073
\[
\psi V=\psi V'=\psi V''=\ldots=\psi V^{(n-1)}.
\]

% ENSEG:W09:PDF068:0074
The value of $F$ can therefore be determined rationally.

% ENSEG:W09:PDF068:0075
2\textsuperscript{o} Conversely, if a function $F$ is rationally determinate and
we set $F=\psi V$, we must have

% ENSEG:W09:PDF068:0076
\[
\psi V=\psi V'=\psi V''=\ldots=\psi V^{(n-1)},
\]

% ENSEG:W09:PDF068:0077
since the equation in $V$ has no commensurable divisor and $V$ satisfies the
equation $F=\psi V$, with $F$ a rational quantity. Thus the function $F$ will
necessarily be invariant under the substitutions of the group written above.

% ENSEG:W09:PDF068:0078
Hence this group has the twofold property stated in the proposed theorem. The
theorem is therefore proved.

% ENSEG:W09:PDF068:0079
We shall call the group in question the \emph{group of the equation}.

% ENSEG:W09:PDF068:0080
\emph{Scholium I.} --- It is clear that, in the group of permutations considered
here, the arrangement of the letters is immaterial; only the substitutions of
letters by which one passes from one permutation to another are to be
considered.

% ENSEG:W09:PDF068:0081
Thus one may choose a first permutation arbitrarily, provided that all the
other permutations are always derived from it by the same substitutions of
letters. The new group so formed will plainly have the same properties as the
first,

\GalSourcePageBoundary{069}{L0068}{40}{40}

% ENSEG:W09:PDF069:0082
since the preceding theorem concerns only the substitutions that can be made
in the functions.

% ENSEG:W09:PDF069:0083
\emph{Scholium II.} --- The substitutions are independent even of the number of
roots.

% ENSEG:W09:PDF069:0084
\begin{center}{\bfseries PROPOSITION II.}\end{center}
\GalSourceErrorMark{W09-SE002}

% ENSEG:W09:PDF069:0085
\noindent\textsc{Theorem}\textsuperscript{(*)}. --- \emph{If one adjoins to a given equation a
root $r$ of an irreducible auxiliary equation, 1\textsuperscript{o} one of two things will
happen: either the group of the equation will remain unchanged, or it will
split into $p$ groups belonging respectively to the given equation when each
root of the auxiliary equation is adjoined; 2\textsuperscript{o} these groups will have the
remarkable property that one passes from one to another by making the same
substitution of letters in every permutation of the first.}

% ENSEG:W09:PDF069:0086
\GalCriticalNote{W09-SE002}{\textbf{Status: repaired by the field-intersection/index theorem.}
The 1897 statement retains $p$ after the phrase defining a prime degree $p$ was
deleted. For original splitting field $L$, $M=K(r)$, $E=L\cap M$, and
$H=\operatorname{Gal}(L/E)$, the correct relation is
$[G:H]=[E:K]$, which divides $[M:K]$. The printed unchanged-or-$p$-equal-groups
dichotomy follows when $[M:K]$ is prime, but not in arbitrary composite degree;
an exact degree-four counterexample is recorded in the certificate. This affects
Propositions II--III and later decompositions. These downstream uses survive
under the stated repair. Tannery 1908 records the defect; Dicker 2026 is later
reconstruction, not historical priority. Evidence: W09-E069, W09-E070,
W09-MATH005, W09-MATH006, and POST-P13-A011.}

% ENSEG:W09:PDF069:0087
\GalCriticalNote{W09-WV003}{\textbf{Status: documented revision history; not a silent repair.}
The 1908 collation records the deleted words ``of prime degree $p$,'' Galois's
note that the proof needed completion, and Liouville's drafted prime-degree
clarification. The 1897 statement above remains translated as printed. This
witness history explains the undefined $p$ and supports, but does not replace,
the field-theoretic repair in W09-SE002. Evidence: W09-E1908-015,
W09-E1908-016, and POST-P13-A014.}

% ENSEG:W09:PDF069:0088
1\textsuperscript{o} If, after $r$ is adjoined, the equation in $V$ considered above
remains irreducible, it is clear that the group of the equation will remain
unchanged. If, on the contrary, it becomes reducible, the equation in $V$ will
then decompose into $p$ factors, all of the same degree and of the form

% ENSEG:W09:PDF069:0089
\[
f(V,r)\times f(V,r')\times f(V,r'')\times\ldots,
\]

% ENSEG:W09:PDF069:0090
where $r,r',r'',\ldots$ are other values of $r$. Thus the group of the given
equation will also decompose into groups, each containing the same number of
permutations, since each value of $V$ corresponds to a permutation. These
groups will be respectively

% ENSEG:W09:PDF069:0091
\GalHistoricalNote{\textsuperscript{(*)} In the statement of the theorem, after the words ``the root $r$ of an
irreducible auxiliary equation,'' Galois first wrote ``of prime degree $p$,''
which he later erased. Likewise, in the proof, instead of ``$r,r',r'',\ldots$
being other values of $r$,'' the original wording read ``$r,r',r'',\ldots$
being the various values of $r$.'' Finally, the following note by the author is
found in the margin of the manuscript:

``There is something to complete in this proof. I do not have time.''

This line was dashed onto the paper very rapidly; that circumstance, together
with the words ``I do not have time,'' makes me think that Galois reread his
Memoir to correct it before going to the field. \hfill \textsc{(A.~Ch.)}}

\GalSourcePageBoundary{070}{L0069}{41}{41}

% ENSEG:W09:PDF070:0092
those of the given equation when $r,r',r'',\ldots$ are adjoined successively.

% ENSEG:W09:PDF070:0093
2\textsuperscript{o} We saw above that all the values of $V$ were rational functions
of one another. Accordingly, suppose that, when $V$ is a root of $f(V,r)=0$,
$F(V)$ is another root; it is clear that, likewise, if $V'$ is a root of
$f(V,r')=0$, then $F(V')$ will be another, for we shall have

% ENSEG:W09:PDF070:0094
\[
f[F(V),r]=\text{a function divisible by }f(V,r).
\]

% ENSEG:W09:PDF070:0095
Therefore (\emph{Lemma I})

% ENSEG:W09:PDF070:0096
\[
f[F(V'),r']=\text{a function divisible by }f(V',r').
\]

% ENSEG:W09:PDF070:0097
This being established, I say that the group relative to $r'$ is obtained by
making everywhere in the group relative to $r$ one and the same substitution
of letters.

% ENSEG:W09:PDF070:0098
Indeed, if, for example, we have

% ENSEG:W09:PDF070:0099
\[
\varphi_{\mu}F(V)=\varphi_{\nu}(V),
\]

% ENSEG:W09:PDF070:0100
then we shall also have (\emph{Lemma I})

% ENSEG:W09:PDF070:0101
\[
\varphi_{\mu}F(V')=\varphi_{\nu}(V').
\]

% ENSEG:W09:PDF070:0102
Thus, to pass from the permutation $[F(V)]$ to the permutation $[F(V')]$, one
must make the same substitution as is made in passing from the permutation
$(V)$ to the permutation $(V')$.

% ENSEG:W09:PDF070:0103
The theorem is therefore proved.

% ENSEG:W09:PDF070:0104
\begin{center}{\bfseries PROPOSITION III.}\end{center}
\GalSourceErrorMark{W09-SE003}

% ENSEG:W09:PDF070:0105
\noindent\textsc{Theorem.} --- \emph{If all the roots of an auxiliary equation are
adjoined to an equation, the groups considered in Theorem II will have in
addition the property that the substitutions are the same in each group.}

% ENSEG:W09:PDF070:0106
The proof will be found\textsuperscript{(*)}.

% ENSEG:W09:PDF070:0107
\GalCriticalNote{W09-SE003}{\textbf{Status: repaired by the normal-intersection proof.}
The revised proposition is printed with only ``The proof will be found.'' If
$M$ is the splitting field of the auxiliary equation and $E=L\cap M$, then
$E/K$ is normal; hence $H=\operatorname{Gal}(L/E)$ is normal in $G$, and the
conjugate groups have the same set of substitutions. This supplies a proof but
does not alter the unproved 1897 text. The result feeds the decompositions in
Proposition V, whose dependent decomposition survives under this repair.
Evidence: W09-E070, W09-E071, W09-E1908-017, and POST-P13-A011.}

% ENSEG:W09:PDF070:0108
\GalCriticalNote{W09-WV004}{\textbf{Status: documented authorial revision and editorial interpolation.}
The 1908 collation identifies the revised statement as a marginal text dated
1832, says that an earlier proof was deleted, and attributes the deciphering
and interpolation of the erased historical note to Liouville. The 1897
statement and apparatus remain distinct from that witness history. Tannery's
1908 collation supplies the explicit provenance. Evidence: W09-E1908-017 and
POST-P13-A014.}

% ENSEG:W09:PDF070:0109
\GalHistoricalNote{\textsuperscript{(*)} In the manuscript, the statement of the theorem just read is in the margin and
replaces another that Galois had written with its proof under the same title,
\emph{Proposition III}. Here is the original text: \textsc{Theorem.} --- If the equa-}

\GalSourcePageBoundary{071}{L0070}{42}{42}

% ENSEG:W09:PDF071:0110
\GalHistoricalNote{tion in $r$ is of the form $r^{p}=A$, and if the $p$th roots of unity are among
the quantities previously adjoined, the $p$ groups considered in Theorem II
will have, in addition, this \GalPrintedError{property}{W09-PE002} that the substitutions of letters by
which one passes from one permutation to another in each group are the same
for all the groups. Indeed, in this case it is immaterial whether one adjoins
this or that value of $r$ to the equation. Consequently, its properties must be
the same after this or that value is adjoined. Thus its group must be the same
with respect to substitutions (Proposition I, Scholium). Hence, etc.

All this is carefully erased; the new statement bears the date 1832 and shows,
by the way it is written, that the author was in an extreme hurry, which
confirms the assertion I made in the preceding Note. \hfill \textsc{(A.~Ch.)}}

% ENSEG:W09:PDF071:0111
\GalCriticalNote{W09-PE002}{\textbf{Status: proved typographical error in the French witness.}
The 1897 historical note prints \emph{proprieté} without the required accent;
the supported French reading is \emph{propriété}. Both necessarily translate
as ``property,'' so the main translation cannot reproduce the accent defect;
the preserved \texttt{W09-PE002} anchor and this note make it visible. The error
has no mathematical downstream effect. No 1846 comparison reading for that
exact token and no explicit prior notice were established. Evidence: W09-E071,
POST-P13-A009, and POST-P13-A014.}

% ENSEG:W09:PDF071:0112
\begin{center}{\bfseries PROPOSITION IV.}\end{center}

% ENSEG:W09:PDF071:0113
\noindent\textsc{Theorem.} --- \emph{If the numerical value of a certain function of
the roots is adjoined to an equation, the group of the equation will be reduced
so that it contains only those permutations under which that function is
invariant.}

% ENSEG:W09:PDF071:0114
Indeed, by Proposition I, every known function must be invariant under the
permutations of the group of the equation.

% ENSEG:W09:PDF071:0115
\begin{center}{\bfseries PROPOSITION V.}\end{center}

% ENSEG:W09:PDF071:0116
\noindent\textsc{Problem.} --- \emph{Under what condition is an equation solvable by
simple radicals?}

% ENSEG:W09:PDF071:0117
I shall first observe that, to solve an equation, its group must successively
be reduced until it contains only one permutation. For when an equation is
solved, every function of its roots is known, even when it is invariant under
no permutation.

% ENSEG:W09:PDF071:0118
This being established, let us seek the condition that the group of an equation
must satisfy in order that it may be reduced in this way by adjoining radical
quantities.

% ENSEG:W09:PDF071:0119
Let us follow the sequence of possible operations in this solution, treating
the extraction of each root of prime degree as a distinct operation.
% GAL1897_W09 stage 2: modern English translation, PDF 072--079 / L0071--L0078.
% Continues Proposition V across the verified PDF 071/072 source boundary.

\GalSourcePageBoundary{072}{L0071}{43}{43}

% ENSEG:W09:PDF072:0120
Adjoin to the equation the first radical extracted in the solution. Two cases
may occur: either adjoining this radical diminishes the group of permutations
of the equation, or, the extraction being merely preparatory, the group remains
the same.

% ENSEG:W09:PDF072:0121
In any event, after some \emph{finite} number of root extractions the group must
be diminished, for otherwise the equation would not be solvable.

% ENSEG:W09:PDF072:0122
If, upon reaching this point, there were several ways to diminish the group of
the given equation by a single root extraction, then, for what we are about to
say, one should consider only a radical of the lowest possible degree among all
the simple radicals whose individual adjunction diminishes the group of the
equation.

% ENSEG:W09:PDF072:0123
Let $p$ therefore be the prime number representing this minimum degree, so that
the group of the equation is diminished by extracting a root of degree $p$.

% ENSEG:W09:PDF072:0124
We may always suppose, at least as far as the group of the equation is
concerned, that a $p$th root of unity, $\alpha$, is among the quantities
previously adjoined to the equation. Since this expression is obtained by
extractions of roots of degree less than $p$, knowing it will not alter the
group of the equation in any way.

% ENSEG:W09:PDF072:0125
Consequently, by Theorems II and III, the group of the equation must decompose
into $p$ groups having, relative to one another, this twofold property:
1\textsuperscript{o} one passes from one to another by one and the same substitution;
2\textsuperscript{o} all contain the same substitutions.

% ENSEG:W09:PDF072:0126
Conversely, I say that if the group of the equation can be divided into $p$
groups having this twofold property, then by one extraction of a $p$th root and
the adjunction of that $p$th root the group of the equation can be reduced to
one of these partial groups.

% ENSEG:W09:PDF072:0127
Indeed, take a function of the roots that is invariant under every substitution
of one partial group and varies under every other substitution. (For this it is
enough to choose a function symmetric in the various values taken, under all
the

\GalSourcePageBoundary{073}{L0072}{44}{44}

% ENSEG:W09:PDF073:0128
permutations of one partial group, by a function invariant under no
substitution.)

% ENSEG:W09:PDF073:0129
\GalCriticalNote{W09-WV005}{\textbf{Status: documented manuscript-topology variant; not an 1897 error.}
The parenthetical construction sentence crossing PDF 072--073 is printed in
the 1897 main text; the 1908 collation identifies it as a marginal note. Its
printed placement is preserved here, and the witness does not authorize moving
or rewriting it. The mathematical argument is unchanged; Tannery 1908 supplies
the explicit collation. Evidence: W09-E1908-017 and POST-P13-A014.}

% ENSEG:W09:PDF073:0130
Let $\theta$ be this function of the roots.

% ENSEG:W09:PDF073:0131
Apply to $\theta$ one of the substitutions of the total group that is not common
to the partial groups. Let $\theta_{1}$ be the result. Apply the same substitution
to $\theta_{1}$ and let $\theta_{2}$ be the result, and so on.

% ENSEG:W09:PDF073:0132
Since $p$ is prime, this sequence cannot stop before the term
$\theta_{p-1}$; after that we shall have

% ENSEG:W09:PDF073:0133
\[
\GalControlReading{\theta_{p}=\theta_{0},\qquad \theta_{p+1}=\theta_{1}}{W09-CD002},
\]

% ENSEG:W09:PDF073:0134
and so on.

% ENSEG:W09:PDF073:0135
This being established, it is clear that the function
\GalSourceErrorMark{W09-SE004}

% ENSEG:W09:PDF073:0136
\[
\bigl(\theta+\alpha\theta_{1}
        +\alpha^{2}\theta_{2}+\ldots
        +\alpha^{p-1}\theta_{p-1}\bigr)^{p}
\]

% ENSEG:W09:PDF073:0137
\GalCriticalNote{W09-SE004}{\textbf{Status: repaired; the radical-tower step survives.}
The printed proof selects the one displayed Fourier--Lagrange coefficient and
adjoins its $p$th root, but that coefficient can vanish even for a nonconstant
orbit. The repair is to select a nonzero nontrivial Fourier coefficient; discrete
Fourier inversion proves that one exists for every nonconstant orbit. This
qualification is needed for the radical reduction that follows and is not
inserted into the translated display. The bounded search found no item-level
prior notice; the Neumann 2013 record alone is not such notice. Evidence:
W09-E073, W09-MATH001, W09-MATH002, W09-MATH004, and POST-P13-A012.}

% ENSEG:W09:PDF073:0138
will be invariant under all permutations of the total group and consequently
will thereby be known.

% ENSEG:W09:PDF073:0139
If the $p$th root of this function is extracted and adjoined to the equation,
then, by Proposition IV, the group of the equation will contain no substitutions
other than those of the partial groups.

% ENSEG:W09:PDF073:0140
Thus the condition above is necessary and sufficient for the group of an
equation to be reducible by a single root extraction.

% ENSEG:W09:PDF073:0141
Adjoin the radical in question to the equation; we can now reason about the new
group as about the preceding one, and it must itself decompose in the indicated
way, and so on, until a group is reached that contains only one permutation.

% ENSEG:W09:PDF073:0142
\emph{Scholium.} --- This procedure is readily observed in the known solution of
general equations of the fourth degree. Indeed, these equations are solved by
means of an equation of the third degree, which itself requires the extraction
of a square root. In the natural order of ideas, one must therefore begin with
this square root. On adjoining this square root to the equation of the fourth
degree, the group of the equation, which contained twenty-four substitutions
in all, decomposes into two groups containing only twelve. Denoting the roots
by $a,b,c,d$,

\GalSourcePageBoundary{074}{L0073}{45}{45}

% ENSEG:W09:PDF074:0143
here is one of these groups:

% ENSEG:W09:PDF074:0144
\[
\begin{array}{ccc}
  abcd, & acdb, & adbc,  \\
  badc, & cabd, & dacb,  \\
  cdab, & dbac, & bcad,  \\
  dcba, & bdca, & cbda.
\end{array}
\]

% ENSEG:W09:PDF074:0145
This group now divides in turn into three groups, as indicated in Theorems II
and III. Thus, after extracting a single radical of the third degree, there
remains simply the group

% ENSEG:W09:PDF074:0146
\[
\begin{array}{l}
  abcd,  \\
  badc,  \\
  cdab,  \\
  dcba:
\end{array}
\]

% ENSEG:W09:PDF074:0147
this group divides again into two groups:

% ENSEG:W09:PDF074:0148
\[
\begin{array}{ll}
  abcd, & cdab,  \\
  badc, & dcba.
\end{array}
\]

% ENSEG:W09:PDF074:0149
Thus, after a single extraction of a square root, there remains

% ENSEG:W09:PDF074:0150
\[
\begin{array}{l}
  abcd,  \\
  badc;
\end{array}
\]

% ENSEG:W09:PDF074:0151
which is finally solved by a single extraction of a square root.

% ENSEG:W09:PDF074:0152
One thus obtains either Descartes's solution or Euler's; for although, after
solving the auxiliary equation of the third degree, Euler extracts three square
roots, it is known that two suffice, since the third follows rationally from
them.

% ENSEG:W09:PDF074:0153
We shall now apply this condition to irreducible equations of prime degree.

\GalSourcePageBoundary{075}{L0074}{46}{46}

% ENSEG:W09:PDF075:0154
\begin{center}
{\large\bfseries Application to irreducible equations of prime degree.}
\end{center}

% ENSEG:W09:PDF075:0155
\begin{center}{\bfseries PROPOSITION VI.}\end{center}
\GalSourceErrorMark{W09-SE005}

% ENSEG:W09:PDF075:0156
\noindent\textsc{Lemma.} --- {\itshape An irreducible equation of prime degree cannot
become reducible by adjoining a radical whose index differs from the degree of
the equation itself.}

% ENSEG:W09:PDF075:0157
For if $r,r',r'',\ldots$ are the various values of the radical and $Fx=0$ is
the given equation, then $Fx$ would have to divide into factors

% ENSEG:W09:PDF075:0158
\[
f(x,r)\times f(x,r')\times\ldots,
\]

% ENSEG:W09:PDF075:0159
all of the same degree, which is impossible unless $f(x,r)$ is of the first
degree in $x$.

% ENSEG:W09:PDF075:0160
Thus an irreducible equation of prime degree cannot become reducible unless
its group is reduced to a single permutation.

% ENSEG:W09:PDF075:0161
\GalCriticalNote{W09-SE005}{\textbf{Status: repaired by explicit Kummer and normality hypotheses.}
The unconditional final inference above fails over $\mathbb{Q}$: after one
adjoins the real cube root of $2$, $x^{3}-2$ becomes reducible while a quadratic
factor and a residual group of order $2$ remain. The argument requires the
inherited roots-of-unity/Kummer normality context. That qualification affects
Propositions VI--VII and is not inserted into the translated 1897 claim. The
unqualified statement is refuted by $x^3-2$; the qualified downstream argument
survives. No prior notice was found in the searched corpus. Evidence: W09-E075,
W09-MATH007, W09-MATH008, and POST-P13-A013.}

% ENSEG:W09:PDF075:0162
\begin{center}{\bfseries PROPOSITION VII.}\end{center}

% ENSEG:W09:PDF075:0163
\noindent\textsc{Problem.} --- {\itshape What is the group of an irreducible equation
of prime degree $n$ that is solvable by radicals?}

% ENSEG:W09:PDF075:0164
\GalControlReading{By the preceding proposition}{W09-CD004}, the smallest possible group before the group
containing only one permutation will contain $n$ permutations. Now,
\GalControlReading{a group of permutations}{W09-CD005} of a prime number $n$ of letters cannot be
reduced to $n$ permutations unless one of these permutations is obtained from
another by a cyclic substitution of order $n$. (\emph{See} M.~Cauchy's Memoir,
\emph{Journal de l'École Polytechnique}, 17th issue.) Thus the
penultimate group will be

% ENSEG:W09:PDF075:0165
\GalCriticalNote{W09-WV006}{\textbf{Status: documented editorial bibliographic expansion; not an 1897 error.}
The 1897 printing gives the full Cauchy citation above; the 1908 collation
reports the manuscript abbreviation ``Journal de l'École, XVII.'' The complete
printed citation is preserved, and no mathematical claim changes. Tannery's
1908 collation supplies the explicit provenance. Evidence: W09-E1908-017 and
POST-P13-A014.}

% ENSEG:W09:PDF075:0166
\[
\text{(G)}\qquad
\left\{
\begin{array}{cccccccc}
x_{0},    &x_{1}, &x_{2}, &x_{3}, &\ldots, &x_{n-3}, &x_{n-2}, &x_{n-1}, \\
x_{1},    &x_{2}, &x_{3}, &x_{4}, &\ldots, &x_{n-2}, &x_{n-1}, &x_{0}, \\
x_{2},    &x_{3}, &\ldots, &\ldots, &\ldots, &x_{n-1}, &x_{0}, &x_{1}, \\
\ldots,  &\ldots, &\ldots, &\ldots, &\ldots, &\ldots, &\ldots, &\ldots, \\
x_{n-1}, &x_{0}, &x_{1}, &\ldots, &\ldots, &x_{n-4}, &x_{n-3}, &x_{n-2}.
\end{array}
\right.
\]

% ENSEG:W09:PDF075:0167
where $x_{0},x_{1},x_{2},\ldots,x_{n-1}$ are the roots.

\GalSourcePageBoundary{076}{L0075}{47}{47}

% ENSEG:W09:PDF076:0168
The group immediately preceding this one in the order of decompositions must
now consist of a certain number of groups all having the same substitutions as
this one. I observe that these substitutions may be expressed as follows (set,
in general, $x_{n}=x_{0}$, $x_{n+1}=x_{1}$, $\ldots$; it is clear that each
substitution of group (G) is obtained by replacing $x_{k}$ everywhere by
$x_{k+c}$, where $c$ is a constant).

% ENSEG:W09:PDF076:0169
Consider any one of the groups similar to group (G). By Theorem II, it must be
obtained by making one and the same substitution everywhere in this group; for
example, by replacing $x_{k}$ everywhere in group (G) by $x_{f(k)}$, where $f$
is a certain function.

% ENSEG:W09:PDF076:0170
Since the substitutions of these new groups must be the same as those of group
(G), we must have

% ENSEG:W09:PDF076:0171
\[
\GalControlReading{f(k+c)=f(k)+C}{W09-CD003},
\]

% ENSEG:W09:PDF076:0172
where $C$ is independent of $k$.

% ENSEG:W09:PDF076:0173
Therefore

% ENSEG:W09:PDF076:0174
\[
\begin{aligned}
f(k+2c)&=f(k)+2C,\\
&\hspace{2em}\dotfill\\
f(k+mc)&=f(k)+mC.
\end{aligned}
\]

% ENSEG:W09:PDF076:0175
If $c=1$ and $k=0$, we find

% ENSEG:W09:PDF076:0176
\[
f(m)=am+b
\]

% ENSEG:W09:PDF076:0177
or

% ENSEG:W09:PDF076:0178
\[
f(k)=ak+b,
\]

% ENSEG:W09:PDF076:0179
where $a$ and $b$ are constants.

% ENSEG:W09:PDF076:0180
Thus the group immediately preceding group (G) can contain only substitutions
such as

% ENSEG:W09:PDF076:0181
\[
x_{k},\qquad x_{ak+b},
\]

% ENSEG:W09:PDF076:0182
and consequently will contain no cyclic substitution other than that of group
(G).

\GalSourcePageBoundary{077}{L0076}{48}{48}

% ENSEG:W09:PDF077:0183
One reasons about this group as about the preceding one, and it follows that the
first group in the order of decompositions---that is, the \emph{actual} group of
the equation---can contain only substitutions of the form

% ENSEG:W09:PDF077:0184
\[
x_{k},\qquad x_{ak+b}.
\]

% ENSEG:W09:PDF077:0185
Thus, {\itshape if an irreducible equation of prime degree is solvable by
radicals, the group of that equation can contain only substitutions of the form

% ENSEG:W09:PDF077:0186
\[
x_{k},\qquad x_{ak+b},
\]

% ENSEG:W09:PDF077:0187
where $a$ and $b$ are constants.}

% ENSEG:W09:PDF077:0188
Conversely, if this condition holds, I say that the equation will be solvable by
radicals. Indeed, consider the functions

% ENSEG:W09:PDF077:0189
\[
\begin{aligned}
&(x_{0}+\alpha x_{1}+\alpha^{2}x_{2}+\ldots+\alpha^{n-1}x_{n-1})^{n}=X_{1},\\
&(x_{0}+\alpha x_{a}+\alpha^{2}x_{2a}+\ldots+\alpha^{n-1}x_{(n-1)a})^{n}=X_{a},\\
&(x_{0}+\alpha x_{a^{2}}+\alpha^{2}x_{2a^{2}}+\ldots+\alpha^{n-1}x_{(n-1)a^{2}})^{n}=X_{a^{2}},\\
&\hspace{12em}\ldots
\end{aligned}
\]

% ENSEG:W09:PDF077:0190
where $\alpha$ is an $n$th root of unity and $a$ is a primitive root modulo $n$.

% ENSEG:W09:PDF077:0191
It is clear that every function invariant under cyclic substitutions of the
quantities $X_{1},X_{a},X_{a^{2}},\ldots$ will in this case be immediately known.
Thus one can find $X_{1},X_{a},X_{a^{2}},\ldots$ by M.~Gauss's method for
binomial equations. Hence, etc.

% ENSEG:W09:PDF077:0192
Therefore, for an irreducible equation of prime degree to be solvable by
radicals, it is \emph{necessary} and \emph{sufficient} that every function invariant
under the substitutions

% ENSEG:W09:PDF077:0193
\[
x_{k},\qquad x_{ak+b}
\]

% ENSEG:W09:PDF077:0194
be rationally known.

% ENSEG:W09:PDF077:0195
Thus the function

% ENSEG:W09:PDF077:0196
\[
(X_{1}-X)(X_{a}-X)(X_{a^{2}}-X)\ldots
\]

% ENSEG:W09:PDF077:0197
must be known for every $X$.

\GalSourcePageBoundary{078}{L0077}{49}{49}

% ENSEG:W09:PDF078:0198
It is therefore \emph{necessary} and \emph{sufficient} that the equation giving this
function of the roots have a rational value for every $X$.

% ENSEG:W09:PDF078:0199
If the given equation has all its coefficients rational, the auxiliary equation
giving this function will likewise have all its coefficients rational, and it
will suffice to determine whether this auxiliary equation of degree
$1\mathbin{.}2\mathbin{.}3\ldots(n-2)$ has a rational root, which one knows how
to do.

% ENSEG:W09:PDF078:0200
This is the method that should be used in practice. But we shall present the
theorem in another form.

% ENSEG:W09:PDF078:0201
\begin{center}{\bfseries PROPOSITION VIII.}\end{center}
\GalSourceErrorMark{W09-SE006}

% ENSEG:W09:PDF078:0202
\noindent\textsc{Theorem.} --- {\itshape For an irreducible equation of prime
degree to be solvable by radicals, it is \emph{necessary} and \emph{sufficient}
that, when any two roots are known, all the others follow rationally from them.}

% ENSEG:W09:PDF078:0203
First, it is necessary, for the substitution

% ENSEG:W09:PDF078:0204
\[
x_{k},\qquad x_{ak+b}
\]

% ENSEG:W09:PDF078:0205
never leaving two letters in the same place, it is clear that, after two roots
are adjoined to the equation, Proposition IV requires its group to reduce to a
single permutation.

% ENSEG:W09:PDF078:0206
\GalCriticalNote{W09-SE006}{\textbf{Status: repaired by adding the critical qualifier ``nonidentity.''}
The identity affine substitution $a=1,b=0$ fixes every letter, contrary to the
unqualified printed sentence; every \emph{nonidentity} affine map fixes at most
one point. Inserting that missing qualifier repairs the necessity argument
without changing the theorem's conclusion, but it is not silently inserted
above. Proposition VIII's two-root criterion survives. No prior notice was
found in the searched corpus. Evidence: W09-E078, W09-MATH009, W09-MATH015,
and POST-P13-A013.}

% ENSEG:W09:PDF078:0207
Second, it is sufficient; for in this case no substitution of the group will
leave two letters in the same places. Consequently, the group will contain at
most $n(n-1)$ permutations. Thus it will contain only one cyclic substitution
(otherwise there would be at least $n^{2}$ permutations). Hence every
substitution of the group, $x_{k},x_{f(k)}$, must satisfy the condition

% ENSEG:W09:PDF078:0208
\[
f(k+c)=f(k)+C.
\]

% ENSEG:W09:PDF078:0209
Hence, etc.

% ENSEG:W09:PDF078:0210
The theorem is therefore proved.

% ENSEG:W09:PDF078:0211
\medskip
\noindent\hfill \textsc{E. G.}\hspace{3em}

\GalSourcePageBoundary{079}{L0078}{50}{50}

% ENSEG:W09:PDF079:0212
\begin{center}{\bfseries EXAMPLE OF THEOREM VII.}\end{center}

% ENSEG:W09:PDF079:0213
Let $n=5$; the group will be as follows:

% ENSEG:W09:PDF079:0214
\[
\GalControlReading{
\begin{array}{l}
  abcde \\
  bcdea \\
  cdeab \\
  deabc \\
  eabcd \\[1.2ex]
  acebd \\
  cebda \\
  ebdac \\
  bdace \\
  daceb \\[1.2ex]
  aedcb \\
  edcba \\
  dcbae \\
  cbaed \\
  baedc \\[1.2ex]
  adbec \\
  dbeca \\
  becad \\
  ecadb \\
  cadbe
\end{array}}{W09-CD006}
\]

% ENSEG:W09:PDF079:0215
\GalArticleEndRule
